% SPDX-License-Identifier: MIT
\chapter{The annotation format}
\label{ch:annotation}

Annotation lives in \cmd{.mytrack} files. One file describes one track: its type,
its data, and how it should look. The format is plain text and designed to be
hand-editable, generated by a script, or exported from a spreadsheet.

\section{A minimal file}

\begin{lstlisting}[language=mytrack]
#makeyourtree-annotation 1

[track]
type = color-strip
title = Biogeographic region

[columns]
names = region

[colors]
Neotropic   = #117733
Afrotropic  = #DDCC77
Palearctic  = #332288

[legend]
title = Region
show  = true

[data]
Homo_sapiens        Palearctic
Pan_troglodytes     Afrotropic
Aotus_trivirgatus   Neotropic
\end{lstlisting}

The first line is the magic marker and is required. Everything after it is a
series of \cmd{[section]} blocks.

\section{Sections}

\begin{center}
\begin{tabular}{@{}lp{0.62\linewidth}@{}}
\toprule
\textbf{Section} & \textbf{Contents} \\
\midrule
\cmd{[track]}   & \cmd{type} (required), \cmd{title}, and \cmd{match}. \\
\cmd{[options]} & Appearance settings for this track type. Chapter~\ref{ch:tracks} lists them per type. \\
\cmd{[columns]} & \cmd{names = a, b, c} --- the column headers, used in legends and axis labels. \\
\cmd{[colors]}  & Explicit colour per category, overriding the automatic palette. \\
\cmd{[legend]}  & \cmd{title}, \cmd{show}, \cmd{order}. \\
\cmd{[data]}    & One row per node: a key, then one value per column. \\
\bottomrule
\end{tabular}
\end{center}

Only \cmd{[track]} with a \cmd{type}, and \cmd{[data]}, are strictly required.

\section{The data section}

Each row begins with a key that identifies a node, followed by one value per
column. Three separators are recognised, tried in this order:

\begin{enumerate}
  \item a tab,
  \item a comma,
  \item a run of two or more spaces.
\end{enumerate}

The last of these means a column-aligned table pasted out of a terminal works
without editing. A single space never separates fields, so taxon names may
contain spaces.

Quoting follows the CSV convention of RFC~4180: wrap a field in double quotes,
and write a literal quote inside it twice.

\subsection{Matching rows to nodes}

By default keys are matched against tip names. \cmd{match = all} in the
\cmd{[track]} section allows internal nodes to match too, which is what
\cmd{clade-range} needs.

If the first data line begins with one of \cmd{id}, \cmd{key}, \cmd{name},
\cmd{node}, \cmd{label}, \cmd{taxon}, \cmd{tip} or \cmd{leaf}, it is treated as a
header row rather than as data.

Rows that match nothing are not silently dropped: they are reported, and the
application's import dialog shows a match report before you commit, so you find
out about a naming mismatch before it becomes a wrong figure.

\subsection{Missing is not zero}

An empty field, or a single \cmd{-}, means \emph{no value}. The track draws a gap
there. This is deliberate and it matters: a zero would be a fabricated
measurement, and a reader cannot tell a fabricated zero from a real one.

\section{Importing a spreadsheet}

You do not have to write \cmd{.mytrack} files by hand. \textbf{Annotate
\textbf{\textrightarrow} Import Annotations} (\keys{Ctrl+I}) reads CSV and
tab-separated files directly. The wizard lets you choose the track type, say which
column holds the node key, and see how many rows matched before anything is
added.

The usual workflow is to keep your data in a spreadsheet, export it as CSV, and
import it. Save the resulting project (\cmd{.mytree}) and the track travels with
it.

\section{Leniency and round-tripping}

Unknown sections and unknown keys are \emph{kept}, not discarded, and recorded as
diagnostics. If a future version writes a key this one does not understand, your
file survives a load-and-save cycle unharmed.

A table written by \app{} and read back is identical to the original. The writer
quotes any field whose bare form would be re-read as something else --- a number,
a \cmd{-}, or a boolean.

The only genuinely fatal error is a file that declares no track type.

\section{Colours}

Colours may be given as \cmd{\#RRGGBB}, \cmd{\#RGB}, \cmd{rgb()} or
\cmd{rgba()}. Bare colour \emph{names} are deliberately not accepted in the
\cmd{[colors]} section, because \cmd{red} is far more often a category label than
a colour request --- accepting it would silently merge a category with its own
swatch.

When no explicit colours are given, categories are assigned from a
colour-vision-safe palette. Twenty-one palettes are available, including the
Paul Tol sets, Okabe--Ito, and the perceptually uniform \cmd{viridis} family.
Prefer these over a hand-picked rainbow: roughly 1 in 12 men has a red--green
colour vision deficiency, and a figure they cannot read is a figure that fails.
